%---------------------------------------------------------------------
\subsection{Thèses soutenues \label{app:phdencours}}
%---------------------------------------------------------------------
% (en précisant les dates de début, dates de fin, taux de co-encadrement et les co­-encadrants, les publications et le devenir des docteurs)

J'indique entre parenthèses mon taux d'encadrement, les noms et affiliations des co-encadrants ainsi que l'école doctorale (ED) de rattachement. Les publications associées à la thèse sont indiquées entre crochets.

\begin{description}
  \item[Frédéric Mortier] 1999-2002 (33\% avec A. Bar-Hen, univ. Lille et C. Baril, GEVES), ED de mathématiques de Lille \cite{MRL05}, chercheur au CIRAD
  \item[Paul Delmar] 2002-2005 (50\% avec J.-J. Daudin, \APT), ED ABIES\footnote{ED Agriculture, alimentation, biologie, environnement et santé, \APT, unversité Paris-Saclay} \cite{DRD05,DRL05}, Roche (pharmacie)
  \item[Tristan Mary-Huard] 2003-2006 (100\%), ED de mathématiques d'Orsay \cite{MDR04,MRD07,MaR09}, chargé de recherche à l'INRAE
  \item[Alain Celisse] 2005-2008 (100\%), ED de mathématiques d'Orsay \cite{CeR08,GRC09,CeR10}, professeur à l'université Panthéon Sorbonne
  \item[Guillem Rigaill] 2007-2010, (50\% avec E. Barillot, Institut Curie et T. Dubois, INSERM), ED ABIES \cite{RLR11,PLH11,VBG13,CKL14}, directeur de recherche à l'INRAE.
  \item[Caroline Bérard] 2008-2011 (50\% avec M.-L. Martin, INRAE), ED ABIES \cite{MMB08,BMT09,BMB11,RAB11,VBM13}, maîtresse de conférences à l'université de Rouen.
  \item[Stevenn Volant] 2008-2012 (50\% avec M.-L. Martin, INRAE), ED ABIES \cite{VMR12,VBM13}, ingénieur de recherche à l'Institut Pasteur.
  \item[Alice Cleynen] 2010-2013 (75\% avec S. Dudoit, IC Berkeley), ED Mathématiques d'Orsay,  bénéficiaire d'une bourse Fullbright, séjour de 6 mois à UC Berkeley (S. Dudoit) \cite{CDR13,CKL14,ClR16,GCR16}, directrice de recherche au CNRS.
  \item[Antoine Channarond] 2010-2013 (75\% avec J.-J. Daudin, \APT), ED Mathématiques d'Orsay \cite{CDR12}, maître de conférences à l'université de Rouen.
  \item[Elisabeta Bonafede] 2012-2015 (20\% avec C. Viroli, univ. Bologne), ED Bologne \cite{BPR15}, ingénieure de recherche en biostatistique dans un groupe pharmaceutique.
  \item[Souhil Chakar] 2012-2015 (50\% avec E. Lebarbier, \APT), EDMH\footnote{ED de Mathématiques Hadamard, université Paris-Saclay.} \cite{CLL17}, professeur agrégé (PRAG) à l'université de Strasbourg.
  \item[Julie Aubert] 2013-2017 (33\% avec S. Schbath, INRA), ED GAO, Evry \cite{ASR21}, ingénieure de recherche à l'INRA.
  \item[Loïc Schwaller] 2013-2016 (75\% avec M. Stumpf, Imperial College), EDMH, 1 séjour de 3 mois à Imperial College (M. Stumpf), 
  \cite{ScR15,ScR17,BSR19,SRS19}, artiste numérique.
  \item[Paul Bastide] 2014-2017 (50\% avec M. Mariadassou, INRA), EDMH, bénéficiaire d'une bourse Fullbright, séjour de 4 mois à Univ. Madison (C. Ané), \cite{BMR17,BAR18,BMR22,BMR24}, 
  chargé de recherche au CNRS.
  \item[Raphaelle Momal] 2017-2020 (50\% avec C. Ambroise, univ. d'Evry), EDMH, \cite{MRA20,MRA21}, ingénieure de recherche à Owkin (biotechnologies et intelligence artificielle).
  \item[Tam Le-Minh] 2020-2023 (33\% avec S. Donnet, INRAE et F. Massol, CNRS), EDMH, 3 articles \cite{LDM25} dont 2 seul (\url{arxiv.org/abs/2103.12597}, \url{arxiv.org/abs/2401.07876}), post-doc à Griffith University, Australie.
\end{description}
 
%---------------------------------------------------------------------
\subsection{Thèses en cours }
%---------------------------------------------------------------------
% (en précisant les dates de début, le taux de co-encadrement et les co-encadrants et les publications)

\begin{description}
  \item[Barbara Bricout] 2022-, (25\% avec S. Donnet, INRAE, T. Galewski, univ. Montpellier / Institut de la Tour du Valat et P. Defos du Rau, Office Français de la Biodiversité), ED Sciences Mathématiques de Paris Centre, 1 article soumis \cite{BDD26}. 
  Cette thèse bénéficie d’un financement de l’Office Français de la Biodiversité (OFB) et de l’Agence Française de Développement (AFD). Un financement d'une année complémentaire a été obtenu de la part de l'Institut des mathématiques pour la planète Terre (IMPT).
  \item[Bixuan Liu] 2024-, (40\% avec V. C. Tran, INRIA et E. Thébault, IEES), ED Sciences Mathématiques de Paris Centre, 1 article soumis seule (\url{arxiv.org/abs/2504.03466}). 
  Cette thèse est financée par la chaire Modélisation Mathématique et Biodiversité (MMB) portée par l’École Polytechnique, le Muséum national d'histoire naturelle et Veolia Environnement.  
  \item[Paul Guillot] 2024-, (33\% avec A. Goddichon, LPSM et L. Sansonnet, LPSM), ED Sciences Mathématiques de Paris Centre, 1 article soumis \cite{GGR26}
\end{description}
